\subsection{Inteligencia artificial y tecnología en la enseñanza de la trigonometría}

La enseñanza de la trigonometría ha incorporado progresivamente herramientas tecnológicas para facilitar la comprensión de conceptos abstractos y promover el pensamiento matemático.

\cite{gonzalezposadaacostaPotencialidadesSoftwareGeogebra2017} investigaron las potencialidades del software GeoGebra en la enseñanza de la trigonometría, demostrando que las representaciones dinámicas y visuales favorecen la comprensión de las funciones trigonométricas. El uso de este tipo de software permite a los estudiantes manipular parámetros y observar en tiempo real los cambios en las gráficas, facilitando el establecimiento de relaciones entre diferentes representaciones matemáticas.

De manera complementaria, \cite{herreracastanedaEnsenanzaConceptosBasicos2013} exploraron la enseñanza de conceptos básicos de trigonometría utilizando tecnología educativa, destacando que los entornos digitales pueden reducir las dificultades asociadas con la abstracción matemática. Los autores encontraron que los estudiantes expuestos a estas metodologías mostraban mayor motivación y comprensión de los conceptos trigonométricos.

Estos trabajos evidencian que la incorporación de tecnología, incluyendo herramientas con capacidades de inteligencia artificial, puede potenciar el aprendizaje de la trigonometría al ofrecer experiencias interactivas y adaptadas a las necesidades individuales de los estudiantes.
